\section{The Category Change: Why 48 Hours Is Not ``3.6 Days Minus 2 Days''}
\label{sec:category}

The central argument of this paper requires a brief excursion into how logistics decisions
are actually made. The economists' instinct is to think of transit time as a continuous
variable: faster is better, and a 50\% improvement is worth 50\% of the benefit of a
100\% improvement. This is approximately correct for pure speed preferences. It is
completely wrong when transit time interacts with physical constraints.

\subsection{Shelf Life as a Binary Constraint}

A California strawberry has a shelf life of five to seven days when picked ripe and held
at 34 degrees Fahrenheit \citep{USDA_NASS2023}. If the berry leaves Watsonville at 6:00 a.m.
Monday, it arrives in New York at approximately 10:00 p.m. Thursday under current conditions ---
roughly 88 hours, or 3.7 days. It has one to three days of shelf life remaining at the
point of arrival. Retailers pay a discount price for a fruit that is already marginal.
Many lots are composted.

This is not a preference problem. A shipper who ``values'' faster transit more highly
cannot get faster transit --- not reliably. The 87th-percentile commitment window for
a solo driver on today's General Purpose interstate lanes is 87 hours \citep{ROUTE_C1}.
That is the number a supply chain planner must build around: not the median, not the
average, but the worst-case commitment time needed to guarantee 95\% on-time delivery.

At 87 hours, the strawberry is marginal. At 48 hours, it is fresh. There is no intermediate
option that unlocks a proportionally intermediate benefit. The economics are binary:
the berry either arrives in time to be valuable, or it doesn't.

The same logic applies across product categories. A biologic drug stable at 2--8 degrees
Celsius for 72 hours \textit{cannot} travel 87 hours by refrigerated truck and arrive
within specification. It can travel 48 hours and arrive with a 24-hour margin. Cold-chain
pharmaceuticals do not fly because airlines are more convenient. They fly because there is
no compliant truck option at current transit times.

\subsection{The Cost Gap That Makes This Matter}

The consequence of this binary constraint is that an enormous volume of freight pays a
15--20$\times$ cost premium --- not because 15--20$\times$ faster service is worth 15--20$\times$
the price, but because no cheaper alternative exists \citep{BTS_airfreight2023}.

\begin{table}[h]
\centering
\caption{Air freight versus refrigerated truck cost by product category (NY--LA corridor)}
\label{tab:cost_comparison}
\begin{tabular}{lrrrl}
\toprule
Product Category & Air (\$/lb) & Truck (\$/lb) & Premium & Min. Transit (truck) \\
\midrule
Biologics / cold-chain drugs  & 8.00--15.00 & 0.80--1.50 & 10--12$\times$ & 48h (2--8\,\textdegree C) \\
Premium produce (strawberries) & 3.00--5.00  & 0.35--0.60 & 8--10$\times$ & 48h (shelf life) \\
Stone fruit / avocados         & 2.50--4.00  & 0.30--0.50 & 8$\times$       & 48h (ripe) \\
Fashion / apparel samples      & 3.00--6.00  & 0.20--0.40 & 15$\times$      & 48h (market window) \\
Electronics components         & 4.00--8.00  & 0.25--0.50 & 16$\times$      & 72h (acceptable) \\
Pacific seafood (live/fresh)   & 5.00--10.00 & 0.40--0.70 & 12$\times$      & 36h (live), 48h (fresh) \\
Industrial parts (emergency)   & 2.00--4.00  & 0.20--0.35 & 10$\times$      & 48h (line-stop cost) \\
\bottomrule
\end{tabular}
\end{table}

\noindent
Table~\ref{tab:cost_comparison} is not a ranking of shipper preferences. It is a catalog of
products that have no choice but to fly today. If the minimum viable truck transit time
is 87 hours and the shelf life (or stability window, or market window) is 72 hours,
the shipper pays the air premium. There is no alternative.

Remove that constraint --- bring the reliable truck transit time below 48 hours ---
and the economics change completely. Not incrementally. Completely.

\subsection{The Simulation Results}

The transit time estimates in this paper derive from a Monte Carlo simulation of 10,000
coast-to-coast trips, implemented in the ROUTE corridor simulation model
\citep{ROUTE_C1}.\footnote{The simulation models weather events, construction delays,
incident frequencies, and Hours of Service compliance using historical FHWA incident data
and NOAA weather records for the I-80 corridor. The model is calibrated against ATRI
speed data for the NY--LA corridor \citep{ATRI_shortage2024}.}

Four scenarios were simulated. Two differ in road infrastructure (current General Purpose
lanes versus I2.0 managed freight lanes). Two differ in driver configuration (solo driver,
team driving, and relay network). Table~\ref{tab:simulation} presents the results.

\begin{table}[h]
\centering
\caption{Monte Carlo simulation results: 10,000 NY--LA trips}
\label{tab:simulation}
\begin{tabular}{llrrrl}
\toprule
Scenario & Infrastructure & p50 (h) & p75 (h) & p95 (h) & 48h SLA \\
\midrule
Solo driver    & GP lanes (today)    & 68.2 & 76.4 & 87.0 & Not achievable \\
Solo driver    & I2.0 managed lanes  & 57.3 & 64.1 & 73.5 & Not achievable \\
Team driving   & I2.0 managed lanes  & 34.1 & 38.6 & 43.8 & 99.5\% of trips \\
Relay network  & I2.0 managed lanes  & 35.6 & 39.8 & 45.4 & 98.8\% of trips \\
\bottomrule
\multicolumn{6}{l}{\footnotesize Source: \citet{ROUTE_C1}. Relay capex: \$40M (8 stations $\times$ \$5M).} \\
\end{tabular}
\end{table}

Two observations from Table~\ref{tab:simulation} are worth dwelling on.

\textit{First}, managed lanes alone do not solve the problem for solo drivers. The p95 drops
from 87 hours to 73.5 hours --- a meaningful improvement for shippers who can tolerate
three-day transit, but not enough to cross the 48-hour threshold. The category change
requires both managed lanes and driver relay.

\textit{Second}, relay and team driving achieve nearly identical transit times (45.4 hours
versus 43.8 hours at p95). But they have very different economics and workforce implications.
Team driving requires two drivers to ride together for 41+ hours, competing for a single bunk
and accumulating fatigue. The relay model assigns each driver one seven-hour leg and sends
them home. The operational cost is similar; the workforce model is radically different.
We return to this in Section~\ref{sec:relay}.

\subsection{The Relay Network Design}

The relay network for the New York--Los Angeles corridor consists of eight stations,
positioned at the major interchange hubs that anchor the I-80 and I-76/I-70 corridor.
Each station is co-located with an existing truck stop or industrial area at or near a
T1-tier diamond interchange \citep{ROUTE_C2}, minimizing new land requirements.

\begin{table}[h]
\centering
\caption{Relay station locations: NY--LA corridor}
\label{tab:relay_stations}
\begin{tabular}{lllrr}
\toprule
Station & Location & Highway Junction & Miles from NY & Leg (mi) \\
\midrule
1 & South Brunswick, NJ     & I-95 / NJ Tpke             &    45 &  45 \\
2 & Cleveland, OH           & I-80 / I-77                &   460 & 415 \\
3 & Chicago, IL             & I-80 / I-90                &   790 & 330 \\
4 & Omaha, NE               & I-80 / I-29                & 1,260 & 470 \\
5 & Salt Lake City, UT      & I-80 / I-15                & 1,980 & 720 \\
6 & Reno, NV                & I-80 / US-395              & 2,400 & 420 \\
7 & Sacramento, CA          & I-80 / I-5                 & 2,630 & 230 \\
8 & Los Angeles (terminal)  & I-10 / I-5                 & 2,790 & 160 \\
\bottomrule
\multicolumn{5}{l}{\footnotesize Average leg: 440 miles. Estimated drive time: 6.7 hours at 65 mph average.} \\
\end{tabular}
\end{table}

The HOS arithmetic is straightforward. Under current Federal Motor Carrier Safety Administration
rules, a solo driver may drive 11 hours within a 14-hour on-duty window, then must take a
mandatory 10-hour rest break \citep{FHWA_freight2023}. A 2,790-mile trip at 60 miles per hour
average requires 46.5 hours of driving. With mandatory rest stops, a solo driver completes
the trip in 46.5 hours of driving plus at least 30 hours of mandatory rest --- a minimum
of 76.5 hours, and more when weather, construction, and incident delay are included.
The relay model eliminates all 30 mandatory rest hours. The truck moves continuously.

\subsection{Why Managed Lanes Matter Even with Relay}

A reader might ask: if the relay network eliminates 30 hours of rest stops, why are
managed lanes necessary? The answer is Donner Pass.

The I-80 corridor crosses the Sierra Nevada at Donner Pass, elevation 7,227 feet.
FHWA records indicate approximately 50 full-closure events per year at this location ---
snowstorms, jackknifed trucks, chain control enforcement \citep{FHWA_freight2023}.
Each closure adds four to twelve hours to transcontinental transit. At 50 closures per year,
this adds roughly six hours to the expected transit time on any given trip (50 closures
$\times$ 8 hours average $\div$ 365 trips per year per lane). But it is not the mean
effect that matters for supply chain planning; it is the tail. A shipper committing to
a 48-hour SLA must plan for the worst 5\% of trips. Donner Pass closures are the primary
driver of the gap between the relay p50 (35.6h) and p95 (45.4h).

The I2.0 managed lane program on the Sierra Nevada segment includes a hardened alignment
that bypasses the most closure-prone section --- either via tunnel or a lower-elevation
alternate alignment. Without this, the relay network still reduces median transit time
dramatically but cannot reliably hold p95 below 48 hours during winter months.
The managed lanes and the relay network are complements, not substitutes.
